\begin{proof}
Multiplying the weight by a constant, assume $w_0=1$. Put
$P_n=\log w_{-n}$ and $Q_n=-\log w_n$ for $n\geq0$.
These sequences are convex, vanish at zero, and grow quadratically.
For either sequence $G=P,Q$, convexity gives
$G_n-G_{n-1}\geq G_n/n$. Put $\beta_j=w_{j+1}/w_j$. For $n\geq1$,
\[
 \beta_{-n}=e^{-(P_n-P_{n-1})}\leq w_{-n}^{-1/n},\qquad
 \beta_{n-1}=e^{-(Q_n-Q_{n-1})}\leq w_n^{1/n},
\]
The signed-root sum therefore gives
\begin{equation}\label{eq:p3-ratio-sum}
 \beta_j\leq1,
 \qquad \sum_{j\in\mathbb Z}\beta_j<\infty.
\end{equation}
In particular, $R$ is bounded. If $|c_j|w_j\leq1$ and
$|b_j|/w_j\leq1$, then for each integer $r\geq1$,
\[
 \sum_j|b_jc_{j-r}|
 \leq\sum_j\frac{w_j}{w_{j-r}}
 \leq\sum_j\beta_{j-r}<\infty.
\]

We prove the following support assertion: if nonzero $b,c$ satisfy
these bounds and
\begin{equation}\label{eq:p3-correlations}
 \sum_jb_jc_{j-r}=0\qquad(r\geq1),
\end{equation}
then $\operatorname{supp}b$ has a finite maximum,
$\operatorname{supp}c$ has a finite minimum, and
\begin{equation}\label{eq:p3-support-separation}
 \max\operatorname{supp}b\leq\min\operatorname{supp}c.
\end{equation}
Independent scaling of $b,c$ gives the same assertion when the two
weighted sequences are merely bounded. We will apply this assertion to
a vector in an invariant subspace and to a continuous annihilator of
that subspace. After proving support separation, the zero-shift pairing
will identify the exact coordinate tail.

We first record an elementary fact about discrete convex envelopes.
Let $f_0=0$ and $f_n\in[0,\infty]$ for $n\geq1$, with $f_n$ finite
at arbitrarily large indices. Suppose there is a finite convex sequence
$G\leq f$ with $G_0=0$ and $G_n\geq\delta n^2$ for all sufficiently
large $n$, where $\delta>0$. The pointwise supremum $F$ of all finite
convex minorants of $f$ anchored at zero is finite: for any $n$, a finite
obstacle at an index $k>n$ bounds every minorant at $n$ by the chord
from zero to $f_k$.
The supremum is convex, since the defining midpoint inequality is
preserved under pointwise suprema. It is therefore the greatest convex
minorant, with $F_0=0$ and $F\geq G$.
If $n\geq1$ and $F_{n+1}-2F_n+F_{n-1}>0$, then
\begin{equation}\label{eq:p3-contact}
 F_n=f_n<\infty.
\end{equation}
Indeed, otherwise raising $F_n$ slightly preserves convexity at $n$,
strengthens it at neighboring indices, and stays below $f$, contradicting
maximality. Also the first differences of $F$ tend to infinity, so
there are arbitrarily large indices satisfying this strict inequality.
We call them contact indices.

First suppose $h=\min\operatorname{supp}c$ exists. If
$\operatorname{supp}b$ were unbounded above, take the greatest convex
minorant $B$ of $f_0=0$, $f_n=-\log|b_n|$ for $n\geq1$, where
$-\log0=\infty$. The obstacle at zero is an artificial anchor; the actual coefficient only
needs $|b_0|\leq1$, which follows from normalization.
The sequence $Q$ is a minorant, so $B\geq Q$.
Writing $v_n=e^{-B_n}$, we have $|b_n|\leq v_n\leq w_n$ for
$n\geq0$, and $|b_p|=v_p$ at arbitrarily large contact indices $p$.
For such $p>h$, the correlation at $r=p-h$ gives
\begin{equation}\label{eq:p3-bounded-side}
 |c_h|\leq\sum_{k\geq1}\frac{v_{p+k}}{v_p}|c_{h+k}|.
\end{equation}
Fix an integer $p_0\geq\max(1,h+1)$. For $p\geq p_0$, convexity
of $B$ and \eqref{eq:p3-ratio-sum} give
\[
 \frac{v_{p+k}}{v_p}|c_{h+k}|
 \leq\frac{v_{p_0+k}}{v_{p_0}w_{h+k}}
 \leq\frac{w_{p_0+k}}{v_{p_0}w_{h+k}}
 \leq\frac{\beta_{h+k}}{v_{p_0}}.
\]
This is summable in $k$. For each fixed $k\geq1$ the left side tends
to zero as $p\to\infty$, because the first differences of $B$ tend
to infinity. Dominated convergence contradicts
\eqref{eq:p3-bounded-side}. Hence $t=\max\operatorname{supp}b$
exists. If $t>h$, the correlation with $r=t-h$ consists of the single
nonzero term $b_tc_h$, a contradiction. Thus $t\leq h$.

The corresponding argument when $b$ is known to be bounded above is
obtained by applying this case to
\[
 \widetilde w_j=1/w_{-j},\qquad
 \widetilde c_j=b_{-j},\qquad
 \widetilde b_j=c_{-j}.
\]
These weights satisfy the same hypotheses, and their correlations are
\eqref{eq:p3-correlations} after reindexing. Thus it remains only to
exclude the case in which $\operatorname{supp}c$ is unbounded below
and $\operatorname{supp}b$ is unbounded above.

In that case let $A_m$ be the convex envelope of
$-\log|c_{-m}|$ for $m\geq1$, anchored at zero, and let $B_p$ be the
convex envelope of $-\log|b_p|$ for $p\geq1$, also anchored at zero.
Both envelopes have value zero at the anchor, while the actual
coefficients satisfy only $|c_0|,|b_0|\leq1$.
Their convex minorants $P,Q$ give $A\geq P$, $B\geq Q$.
Define
\[
 u_{-m}=e^{A_m},\qquad u_p=e^{-B_p}\qquad(m,p\geq0).
\]
Since $A_0=B_0=0$ and $A_1\geq P_1\geq0$, $B_1\geq Q_1\geq0$,
convexity makes all first differences of $A$ and $B$ nonnegative.
The sequence $u$ is therefore positive and decreasing, with $u_0=1$; its
logarithm is convex on the negative half-line and concave on the
positive half-line. The obstacle inequalities and
$u_j\geq w_j$ for $j\leq0$, $u_j\leq w_j$ for $j\geq0$ give
\begin{equation}\label{eq:p3-envelope-bounds}
 |c_j|\leq u_j^{-1},\qquad |b_j|\leq u_j
 \qquad(j\in\mathbb Z).
\end{equation}
The quadratic lower bounds for $A,B$ also give
\begin{equation}\label{eq:p3-envelope-ratios}
 \sum_{j\in\mathbb Z}\frac{u_{j+1}}{u_j}<\infty.
\end{equation}

Choose $\log3<\lambda<a$. There are arbitrarily large $m$ with
\begin{equation}\label{eq:p3-large-curvature}
 A_{m+1}-2A_m+A_{m-1}>2\lambda.
\end{equation}
Otherwise double summation gives $A_m\leq\lambda m^2+O(m)$,
contrary to $A_m\geq P_m$ and $P_m/m^2\to a$.
Set $\alpha_m=A_m-A_{m-1}$ and $\gamma_p=B_p-B_{p-1}$.
Convexity and the quadratic lower bounds imply that both sequences
are nondecreasing and tend to infinity.
We match these slopes to choose two contact coefficients whose product
dominates the other terms in the middle of one correlation. The two
remaining tails will then be negligible relative to that product.
For each sufficiently large such $m$, choose the last index $p\geq1$ with
\begin{equation}\label{eq:p3-slope-choice}
 \gamma_p\leq\alpha_m+\lambda<\gamma_{p+1}.
\end{equation}
The choice exists because the nondecreasing sequence $\gamma_p$
tends to infinity; moreover $p\to\infty$ with $m$. Equations
\eqref{eq:p3-large-curvature}--\eqref{eq:p3-slope-choice} yield
\[
 \gamma_{p+1}>\alpha_m+\lambda,
 \qquad \alpha_{m+1}>\gamma_p+\lambda.
\]
Both $m$ and $p$ are contact indices. Therefore
\begin{equation}\label{eq:p3-contact-product}
 |b_pc_{-m}|=u_p/u_{-m}=:M>0.
\end{equation}
The sequence
\[
 F(n)=\log\frac{u_{p+n}}{u_{-m+n}}
 =-B_{p+n}-A_{m-n},\qquad -p\leq n\leq m,
\]
is concave. The preceding slope inequalities give
$F(1)-F(0)<-\lambda$ and $F(-1)-F(0)<-\lambda$.
Consequently
\begin{equation}\label{eq:p3-geometric-middle}
 \frac{u_{p+n}}{u_{-m+n}}\leq M e^{-\lambda|n|}
 \qquad(-p\leq n\leq m).
\end{equation}

Put $q=m+p$ and use the correlation for $r=q$, with $j=p+n$.
After isolating the term $n=0$, the bounds
\eqref{eq:p3-envelope-bounds}--\eqref{eq:p3-contact-product} give
\[
 M\leq\sum_{n\ne0}\frac{u_{p+n}}{u_{-m+n}}.
\]
The part with $-p\leq n\leq m$ is at most
\[
 M C_\lambda,\qquad
 C_\lambda=2\sum_{n\geq1}e^{-\lambda n}
 =\frac{2}{e^\lambda-1}<1.
\]
At $n=m$, \eqref{eq:p3-geometric-middle} gives $u_q\leq M$.
Thus the part with $n>m$, divided by $M$, is at most
\[
 \sum_{k\geq1}\frac{u_{k+q}}{u_ku_q}.
\]
Concavity of $\log u$ on the positive half-line implies, for $q\geq1$,
\[
 \frac{u_{k+q}}{u_ku_q}
 \leq\frac{u_{k+1}}{u_1u_k}.
\]
The majorant is summable by \eqref{eq:p3-envelope-ratios}; for each
fixed $k\geq1$, the term on the left tends to zero as $q\to\infty$,
because $u_{q+k}/u_q$ is a product of $k$ step ratios tending to zero.
Dominated convergence shows that this tail is $o(M)$.
For the other tail, put $v_k=1/u_{-k}$. Then $\log v$ is concave,
its step ratios are summable, and
\eqref{eq:p3-geometric-middle} at $n=-p$ gives $v_q\leq M$.
Writing $k=-p-n$ bounds this tail, divided by $M$, by
\[
 \sum_{k\geq1}\frac{v_{k+q}}{v_kv_q}=o(1)
\]
by the same argument. We obtain $1\leq C_\lambda+o(1)$, a
contradiction. This proves \eqref{eq:p3-support-separation}.

Finally let $L$ be a nonzero proper closed $R$-invariant subspace.
Its continuous annihilators have the form
$c\mapsto\sum_jb_jc_j$ with $b/w\in\ell^2$.
For $c\in L$ and such an annihilator $b$, the correlations vanish for
all $r\geq0$, and the two required weighted boundedness conditions
hold. Fix a nonzero annihilator $b^*$ and put
$t=\max\operatorname{supp}b^*$. Support separation shows that every
nonzero $c\in L$ is supported at indices at least $t$; orthogonality
at $r=0$ then forces $c_t=0$. Since $L\ne\{0\}$, the union of supports
of its elements is a nonempty subset of $\{t+1,t+2,\ldots\}$ and
therefore has a least index $h$. Choose $c^*\in L$ with $c^*_h\ne0$.
Every $c\in L$ vanishes below $h$. Conversely, support separation
with $c^*$ shows that every nonzero annihilator $b$ is supported at
indices at most $h$, and orthogonality at $r=0$ forces $b_h=0$.
All annihilators therefore vanish on the coordinate tail beginning
at $h$. Hahn--Banach and closedness imply that this tail is contained
in $L$. Thus $L$ is exactly that tail.

Finally the map $c\mapsto(w_jc_j)_j$ is a unitary from
$\ell^2(w)$ onto $K$. It carries $R$ to the bounded weighted shift
$Te_j=(w_{j+1}/w_j)e_{j+1}$ and preserves every coordinate-vanishing
condition. This proves the last assertion as well.
\end{proof}
